%-------------------------------------------------------------------------------
%  CẤU HÌNH DÙNG CHUNG cho CV (main.tex) và Đơn xin việc (don-xin-viec.tex)
%
%  Sửa thông tin liên hệ ở đây MỘT LẦN là cả hai tài liệu cùng cập nhật.
%-------------------------------------------------------------------------------

\geometry{left=1.4cm, top=.8cm, right=1.4cm, bottom=1.4cm, footskip=.5cm}

% Màu chủ đạo. Các lựa chọn có sẵn của Awesome-CV: awesome-emerald, awesome-skyblue,
% awesome-red, awesome-pink, awesome-orange, awesome-nephritis, awesome-concrete,
% awesome-darknight. Ở đây dùng màu xanh y tế cho phù hợp ngành Dược.
\definecolor{awesome}{HTML}{00796B}

%-------------------------------------------------------------------------------
% SỬA MẪU GỐC CHO HỢP VỚI TIẾNG VIỆT
%-------------------------------------------------------------------------------
% Mẫu gốc tô màu 3 KÝ TỰ đầu của tiêu đề mục (\StrSplit{...}{3}). Với tiếng Việt,
% cách này cắt ngang giữa từ ("Kin|h nghiệm", "Chứ|ng chỉ") trông như lỗi hiển thị,
% nên ở đây tô màu toàn bộ tiêu đề — cho cả CV lẫn đơn xin việc.
\setbool{acvSectionColorHighlight}{false}
\renewcommand*{\sectionstyle}[1]{{\color{awesome}\sectionstyleface{#1}}}
\renewcommand*{\lettersectionstyle}[1]{{\color{awesome}\lettersectionstyleface{#1}}}

% Mẫu gốc dùng chữ small-caps (\scshape) cho dòng chức danh, chân trang, tên vị trí
% và địa chỉ người nhận. Font không có glyph small-caps cho các chữ tiếng Việt có dấu
% (ư, ĩ, ẳ, ơ...) nên chúng rơi về chữ thường, làm dòng chữ cao thấp không đều.
% Thay bằng chữ IN HOA để giữ đúng ý đồ thiết kế mà vẫn hiển thị đủ dấu.
\renewcommand*{\headerpositionstyle}[1]{{\fontsize{7.6pt}{1em}\bodyfont\color{awesome}\MakeUppercase{#1}}}
\renewcommand*{\footerstyle}[1]{{\fontsize{8pt}{1em}\footerfont\color{lighttext}\MakeUppercase{#1}}}
\renewcommand*{\entrypositionstyle}[1]{{\fontsize{8pt}{1em}\bodyfont\color{graytext}\MakeUppercase{#1}}}
\renewcommand*{\subentrypositionstyle}[1]{{\fontsize{7pt}{1em}\bodyfont\color{graytext}\MakeUppercase{#1}}}
% Địa chỉ người nhận giữ nguyên chữ thường cho dễ đọc (tên công ty, đường, phường)
\renewcommand*{\recipientaddressstyle}[1]{{\fontsize{9pt}{1em}\bodyfont\color{graytext} #1}}

% Thu gọn khoảng cách giữa các mục để tài liệu vừa đúng 1 trang
\renewcommand{\acvSectionTopSkip}{2.4mm}
\renewcommand{\acvSectionContentTopSkip}{2mm}

\renewcommand{\acvHeaderSocialSep}{\quad\textbar\quad}

% Ngày tháng theo định dạng Việt Nam (\today của LaTeX trả về tiếng Anh)
\makeatletter
\newcommand{\ngayhomnay}{\two@digits{\day}/\two@digits{\month}/\the\year}
\makeatother

%-------------------------------------------------------------------------------
% THÔNG TIN CÁ NHÂN
%-------------------------------------------------------------------------------
\photo[circle,edge,left]{./avatar}
\name{Võ Vương}{Khánh Vy}
\position{Dược sĩ Cao đẳng~~·~~Dược sĩ tư vấn}
\address{KV. Phụng Thạnh 1, P. Thuận Hưng, TP. Cần Thơ}

\mobile{0949 507 278}
\email{vyk5045@gmail.com}
\dateofbirth{20/10/2005}

\quote{“Tư vấn đúng thuốc, đúng liều, đúng người --- để mỗi khách hàng quay lại vì tin tưởng.”}
